\section{Defining always-on agents}\label{sec:definition}


\regime{Substance}{what an always-on agent is, and the six axes that characterize a unit of its state.}
An always-on agent is best understood not as a memory-augmented language model but as a \emph{persistent-state system}: a system whose behavior at any moment depends on state it has accumulated, transformed, and retained across earlier interactions, and whose correctness depends as much on governing that state as on recalling it. A retrieval cache asks what the system can recall. A persistent-state system also asks whether a recalled item is still current, whether it is authorized to influence the action under consideration, whether it is scoped to the user and task at hand, whether it descends from a trustworthy source, whether the decisions it licenses can be reversed, and whether it is a passive fact or an executable commitment. The definition below separates always-on agents from five adjacent system classes routinely conflated with them, introduces the six state axes that form the survey's analytic vocabulary, traces their cognitive-architecture lineage, and states what the definition includes and excludes, in particular why personalization and trigger behavior are persistent-state \emph{types}.

\subsection{An operational definition}\label{subsec:operational}

We define an always-on agent as an agent whose policy at time $t$ depends on state accumulated before $t$, beyond the current prompt or task instance. The defining property is persistence, not literal continuous execution: an agent that runs once a day, wakes on an external trigger, or is invoked only when a user opens a session is still always-on if its action at invocation is conditioned on durable state carried across invocations. Conversely, an agent that runs continuously but resets all state between requests is episodic, not always-on, because nothing it retains can authorize a later action.

The definition is operational so that it can be checked against a system rather than asserted of it. A system is always-on if it satisfies three conditions. First, it has \emph{persistent identity} across sessions or restarts, so the same agent instance is recoverable after interruption and can be held accountable for what it did before. Second, it maintains \emph{agent-owned, mutable, durable state} about users, tasks, tools, environments, policies, or other agents, state that outlives the interaction in which it was created. Third, it can \emph{use that state for later action}, closing the loop from accumulation back to behavior; state written but never read is inert and does not make an agent always-on. If that state is consequential, it creates a fourth requirement, \emph{temporal accountability}: the system should be able to record what it knew, when it knew it, why it acted, and how state that later proves contaminated can be identified and repaired. The first three properties classify the system. The fourth is the obligation current systems often fail to satisfy, and that failure is the empirical focus of this survey.

Temporal accountability cannot be reduced to retrieval quality. A system can retrieve a fact with perfect recall and still be unable to say whether the fact remains authorized, whether its scope still covers the present request, or which past actions would have to be undone if the fact turns out to be poisoned. The classical agent surveys that mapped the design space of LLM agents organized it around perception, reasoning, memory, and action as functional modules \citep{xi2023riseandpotential,wang2023surveylargelanguagemodelagents}, a useful decomposition that nonetheless treats memory as one black-box module to be written and read, and does not ask what governs the transition from a retrieved item to an authorized influence on the next action. For an always-on agent, that transition is the object of study. An episodic agent can be evaluated on whether its final answer is correct; an always-on agent must be evaluated on whether each state transition was legitimate, because an illegitimate transition can produce a correct-looking answer today and an irreversible harm tomorrow. Later sections therefore evaluate state mutations and action consequences rather than final answers alone.

\paragraph{An objection: is the governance gap rational triage?} A careful skeptic can grant every corpus statistic in this survey and still resist its conclusion. Perhaps rollback is rare because most agent state is intentionally disposable: a wrong preference or a stale cached plan can be overwritten rather than reverted. Perhaps authority is the rarest axis because many deployed agents run in bounded single-user, high-trust contexts where fine-grained authorization would add cost without reducing risk. On this reading the sparse return arc is specialization, not a gap. The claim we make is conditional: \emph{whenever persistent state is both retained and consequential, the absence of governance is an undefended design choice rather than a rational one}. Three lines of evidence make the antecedent common rather than rare. First, production-scale systems accumulate long-horizon consequential state: credential and permission tokens, durable user preferences, learned policies, and task ledgers whose corruption persists across sessions (Sections~\ref{sec:statetax} and~\ref{sec:substrates}). Second, the failure modes traced to governance omissions are documented in attack and incident literature and arise even in the bounded settings the objection invokes (Section~\ref{sec:failures}). Third, the budget-matched and continual-learning results we survey show systems \emph{explicitly designed to learn} degrading below a no-memory baseline over long horizons (Section~\ref{sec:mechanisms}); that failure arose because the governance condition was present and ungoverned. Where state is genuinely ephemeral, this survey's prescriptions do not apply. The burden is to govern the state that is not.

\paragraph{A note on temporal heterogeneity.} The definition spans three execution patterns: continuously running agents, periodically scheduled ones, and session-triggered ones that run on demand but carry state across sessions. We group them because all three retain persistent state and inherit its governance obligations, yet their risk surfaces differ. A continuously running or scheduled agent faces \emph{prospective} governance, when it is appropriate to act on latent state at all, which is the trigger-calibration axis. A session-triggered agent faces \emph{retrospective} governance, whether state carried from a prior session still authorizes the action it is about to take. The corpus coding aggregates the two, so a count such as the twenty-seven works exposing rollback mixes rollback for periodic re-execution with rollback for session recovery, which are mechanically distinct. We flag the conflation so that later empirical frequencies are read with the right grain: the governance axes apply across all three patterns, but the prevalence of any single failure mode may be concentrated in one.

\subsection{Boundaries: five adjacent system classes}\label{subsec:boundaries}

The definition earns its weight by excluding things, so we sharpen it by contrasting always-on agents against five system classes with which they are habitually conflated, because long context, retrieval, context engineering, episodic tool use, and conversational memory all touch information that persists in some sense. The distinction in every case is whether the system governs what becomes durable, authoritative, agent-owned state, or only processes information that happens to be present.

\paragraph{Episodic LLM agents.} The baseline contrast is the episodic tool-using agent, the dominant evaluation paradigm. Frameworks in the ReAct lineage interleave reasoning and action over many steps, binding intermediate reasoning to tool calls and validating outcomes within a task \citep{yao2023reactsynergizingreasoningacting}. Such an agent can act with great sophistication over a long horizon, yet it resets between tasks: nothing it learned in one episode authorizes an action in the next, and the always-on agent breaks that reset. What separates the two is not the length or complexity of a single episode but whether identity, commitments, permissions, and lessons survive the episode boundary. An episodic agent satisfies the survey's governance invariants vacuously, by forgetting everything, which makes it safe in a narrow sense and useless for the longitudinal work always-on agents are built for. The gap is foundational: the episodic paradigm has no vocabulary for an action that is wrong only because of state inherited from a prior session, so its benchmarks cannot detect that failure class at all.

\paragraph{Memory-augmented agents.} The closest neighbor is the memory-augmented agent, which adds an external store, a vector index, a note graph, or an operating-system-style paged context, that the policy writes to and reads from across sessions. MemGPT, for instance, treats memory as an OS-managed hierarchy with explicit paging between an in-context working set and externalized long-term storage, decoupling the effective horizon from the context budget \citep{packer2024memgptllmsoperatingsystems}. This is genuine cross-session state, so on our definition such systems are a subclass of always-on agents; the distinction is internal rather than exclusionary. Memory-augmented agents typically instantiate accumulation and recall but not temporal accountability: they expose write, organize, and retrieve, and leave authority, deletion propagation, and rollback implicit, so that a deletion edits the primary store while summaries, embeddings, and promoted tiers retain the deleted content. The boundary is therefore not whether the system has memory but whether the memory is \emph{governed}; persistent state, as we use the term, is wider than retrievable memory and includes task ledgers, permissions, credentials, commitments, provenance, and externally committed effects that remain linked to the records that authorized them.

\paragraph{Long-context models.} A long-context model can ingest a very large input, an entire codebase, a year of chat history, but it does not decide what to preserve for future use. Long-context benchmarks measure whether a model can use what is already inside its window, and repeatedly show the answer is unreliable even when the information is present: models underuse mid-context information despite attending to its boundaries \citep{liu2023lostmiddlelanguagemodels}, and broad bilingual long-context suites confirm that task performance degrades with length in ways the architecture does not advertise \citep{bai2024longbenchbilingualmultitaskbenchmark}. These results matter for always-on agents as \emph{constraints}, since whatever an agent retrieves must still be usable once it enters the window, but they answer a different question. A long-context model is given its context; an always-on agent must decide which signals become durable state, when they become authoritative, and how they are later revised or revoked. The window is volatile working memory that vanishes at the end of the call; persistent state is what the agent chose to keep, and the keeping is the governed act long context never performs.

\paragraph{RAG systems.} Retrieval-augmented generation retrieves from an external corpus to ground generation, and RAG benchmarks supply the vocabulary of faithfulness, noise robustness, evidence selection, and rejection of irrelevant passages \citep{chen2023benchmarkinglargelanguagemodels}. A RAG system reads from a corpus but does not, in general, write back to it from experience: the corpus is a fixed, externally curated knowledge base, not state the agent accumulates and must later govern. The always-on agent inverts the asymmetry. It populates its store from its own observations and actions, so it inherits problems a static corpus does not have: self-poisoning through an untrusted write, staleness when a once-true fact changes, and the obligation to delete on request and propagate that deletion to every derived copy. RAG faithfulness asks whether retrieved evidence was used correctly; the always-on question is whether the evidence should have been in the store at all, whether its authority is current, and whether it can be removed. RAG mechanics are thus a substrate the always-on agent often uses, but RAG evaluation does not test the write-side governance obligations that define the always-on setting.

\paragraph{Personalized and proactive assistants.} The final and most subtle boundary is the personalized or proactive assistant. A long line of work models users and conditions outputs on user profiles: the LaMP suite evaluates profile-conditioned generation \citep{salemi-etal-2024-lamp}, recent surveys organize the foundations and evaluation of personalized LLM agents \citep{xu2026personalizedllmpoweredagentsfoundations}, PrefEval tests whether a model infers and follows stated preferences across long multi-session context \citep{zhao2025prefeval}, and PersonaMem tracks an evolving user persona over interactions \citep{jiang2025personamem}. Proactivity, deciding when to surface help or act unprompted, is a parallel capability evaluated by datasets such as proactive-agent and proactive-recall benchmarks \citep{lu2024proactiveagent,he2024madialbench}. These systems unambiguously build persistent state, and many are always-on agents. But personalization and proactivity are \emph{capabilities expressed over} persistent state, not its definition. A model can top a preference-following benchmark yet have no mechanism to revoke a preference the user later retracts; high recall of a preference whose authority has lapsed is not a success but an authority-monotonicity failure the benchmark has no field to score. We therefore treat personalization and trigger behavior as state types and use cases, returning to this in Section~\ref{subsec:include-exclude}.

Table~\ref{tab:boundary} consolidates these contrasts along the separating question: does the system govern durable, authoritative, agent-owned state?

\begin{table}[t]
\centering
\small
\caption{Always-on agents against five adjacent system classes. The separating question is not whether information persists but whether the system governs durable, agent-owned state: what becomes authoritative, and how it can be revised, scoped, revoked, and rolled back.}
\label{tab:boundary}
\setlength{\tabcolsep}{4pt}\renewcommand{\arraystretch}{1.2}
\begin{tabularx}{\linewidth}{p{0.18\linewidth}YY}
\toprule
System class & Question it answers & Governs durable state? \\
\midrule
Episodic LLM agent & Can it act correctly within one task before reset? & No: resets between tasks \\
Long-context model & Can it use what is already in the window? & No: does not choose what to keep \\
RAG system & Is external evidence selected and used faithfully? & No: reads a fixed corpus, rarely writes back \\
Memory-augmented agent & Can it recall cross-session facts? & Partly: writes and reads, but governance left implicit \\
Personalized / proactive assistant & Does output match the user and arrive at the right time? & Partly: builds state but rarely scopes, revokes, or rolls back \\
\midrule
Always-on agent & Which retained state is authoritative for the next action, and how is it revised, scoped, revoked, and repaired? & Yes, by definition \\
\bottomrule
\end{tabularx}
\end{table}

The common thread is that each adjacent class makes one part of the always-on problem easy by assuming the rest away. Episodic agents assume away persistence; long-context models assume the context is given; RAG assumes the corpus is curated and read-only; memory-augmented systems assume recall is the whole job; personalized assistants assume that matching the user is the whole job. The always-on agent can assume none of these, and the open problem the boundary analysis exposes is that we lack a single benchmark crossing all of them at once, stressing persistent identity, governed writes, scope, revocation, and rollback in one protocol rather than testing each adjacent capability in isolation.

\begin{table}[t]
\centering
\scriptsize
\setlength{\tabcolsep}{3pt}\renewcommand{\arraystretch}{1.05}
\begin{tabularx}{\linewidth}{p{0.30\linewidth}cY}
\toprule
System example & In scope? & Boundary reason and resulting obligation \\
\midrule
Long-context QA over supplied history & No & Context is supplied for one inference; ordinary context privacy and evidence use apply. \\
Read-only RAG over a curated corpus & No & The agent reads an external corpus but does not own or mutate it; corpus provenance and retrieval faithfulness matter. \\
Episodic ReAct or AppWorld task & No & Identity and state reset at the task boundary; evaluation is task-local. \\
Saved conversational memory & Yes & User facts persist and shape later responses; revocation, scope, provenance, and deletion propagation become obligations. \\
MemGPT/Letta-style external memory & Yes & The agent writes and retrieves durable memory; authority, mutability, and recoverability must be represented. \\
Repository-memory coding agent & Yes & Project conventions and prior fixes affect later edits; provenance, conflict handling, and rollback become relevant. \\
Calendar or trigger agent & Yes & Stored conditions can cause later proactive action; current authority, temporal validity, and confirmation policy matter. \\
Shared workspace or credentialed tool agent & Yes & Multiple principals or durable grants affect later side effects; ownership, permission epochs, idempotency, and compensation are required. \\
\bottomrule
\end{tabularx}
\caption{Decision procedure for the always-on boundary. A system is in scope when it has persistent identity, owns mutable durable state, and can use that state for later action. Temporal accountability is not required to classify the system; it is the governance obligation created when the state is consequential.}
\label{tab:decision}
\end{table}

\subsection{The six state axes}\label{subsec:axes}

If persistent state is the unit of analysis, we need a vocabulary precise enough that a benchmark can perturb an individual state item and a mechanism can be said to govern it. We characterize a persistent state item along six diagnostic axes, each phrased as a question an evaluation can ask of the item and that a system either does or does not answer. The axes are not orthogonal variables. Provenance often supplies evidence for authority; scope constrains authority; actionability determines how strong the authority and recoverability requirements must be; and recoverability depends on provenance being preserved. The axes are still useful because each names a dimension along which a distinct, persistence-specific failure class arises, grounded in either the cognitive-science literature that first distinguished it or the agent literature that exposed its failure mode. Later sections measure how unevenly the field covers them, with authority the rarest at 72 of 435 coded works and mutability the most common at 160.

\paragraph{Authority.} Authority asks who or what permits this state to influence an action, distinguishing a fact the agent only holds from one it is licensed to act upon. The cognitive lineage supplies only a partial analogy: ACT-R separates declarative chunks, which are simply available, from procedural production rules selected by a utility that governs when they fire \citep{anderson2004actr}, prefiguring the distinction between possessing a record and acting on it. That utility is an action-selection preference, not an authorization, so the authority axis in our sense is grounded less in cognitive architecture than in capability-based security and access control, where the right to influence an action is an explicit, revocable grant rather than a learned preference (we develop this lineage in the governance part). In the agent setting, authority is the axis most directly attacked: work on long-term state poisoning shows routine interactions can gradually weaken an agent's confirmation boundaries and expand the scope of actions it will take, an authority drift no recall metric registers \citep{xu2026toxicchats}. The corresponding invariant, authority monotonicity, holds that a record may influence an action only under a current, unrevoked authority, so an expired permission cannot license a write or a tool call. The gap this axis exposes is that the bulk of the literature treats every stored item as equally entitled to act, with no representation of who authorized it or whether that authorization still holds; authority is the least represented axis in our coding.

\paragraph{Scope.} Scope asks which user, task, tool, time window, or group an item may be used for. Its cognitive grounding is the long-standing distinction between volatile working memory bound to the current situation and durable long-term memory, formalized in the multi-component working-memory model whose episodic buffer binds integrated representations for the task at hand \citep{baddeley2000episodic}. Scope keeps one user's preference from acting on behalf of another, and it is what a category-bounded memory enforces when an in-car assistant restricts each stored preference to its declared category and prunes redundant or contradictory entries \citep{kirmayr2025carmem}. The invariant is scope non-expansion: a transition may narrow or preserve scope but never silently widen it. The gap is that even systems enforcing scope tend to do so statically, fixing categories or roles at design time, with no account of what happens when a scope must be redefined, when authority changes dynamically, or when an item legitimately moves between scopes, leaving cross-scope leakage as a standing risk.

\paragraph{Mutability.} Mutability asks whether an item can be revised, superseded, decayed, or locked, and on what timescale. This axis is grounded in the quantitative study of forgetting, from the reproducible replication of the Ebbinghaus forgetting curve with fitted decay functions \citep{murre2015replication} to the complementary-learning-systems account in which fast and slow stores change at deliberately different rates \citep{mcclelland1995complementary}. In the agent setting, mutability is exposed by benchmarks showing user attributes, habits, and preferences drift on different timelines driven by external events, so a single global update rule is wrong \citep{xie2026dynamicmem}, and by principled-forgetting benchmarks balancing recall against time-decay invalidation \citep{uddin2026recallforgettingbenchmarkinglongterm}. Mutability is the best-covered axis because forgetting and updating are the most studied operations, yet the coverage is shallow: most work decays uniformly, whereas real state is heterogeneous, some items permanent, others ephemeral, and the open problem is authority-conditioned mutability where the right to revise or decay an item is itself governed.

\paragraph{Provenance.} Provenance asks which source, timestamp, and chain of transformations produced an item. Its cognitive root is the episodic-versus-semantic distinction: episodic memory is memory for events tied to a time and place, while semantic memory abstracts away the source \citep{tulving1972episodic}, and provenance is the agent-side discipline of not losing the episodic source when an item is consolidated into a semantic fact. The complementary-learning-systems literature makes the stakes explicit, since consolidation through replay is the operation that can preserve or discard the trace of origin \citep{mcclelland1995complementary,wilson1994reactivation}. The invariant is provenance preservation: consolidation may compress an item's value but must retain enough of its derivation to verify, attribute, and revise the result later. The exposed gap is that the dominant consolidation methods are lossy by design and flatten provenance, so an agent loses the ability to say where a fact came from at precisely the moment it most needs that, when the fact is later challenged or found poisoned.

\paragraph{Recoverability.} Recoverability asks whether the derived state and the decisions an item caused can be rolled back. This is the axis the field covers least well after authority, and it has the weakest cognitive analogue, since biological forgetting is largely irreversible, which is itself instructive: the canonical forgetting curve describes loss without recovery \citep{murre2015replication}, and ACT-R's activation decay offers no undo for a chunk fallen below threshold \citep{anderson2004actr}. Artificial agents are not bound by this limit, but most inherit it by default. Recoverability is what the principled-forgetting line gestures at when it distinguishes accessibility decay from true deletion \citep{uddin2026recallforgettingbenchmarkinglongterm}, yet even there recovering the downstream actions a forgotten fact licensed is left open. The invariant is rollback traceability: every action carries a handle back to the records that justified it, so a later-discovered bad record yields a bounded set of decisions to repair. The gap is stark and quantified later: rollback is the rarest lifecycle stage in the coded corpus, and recoverability is one place where the survey's governance framing does real organizing work.

\paragraph{Actionability.} Actionability asks what kind of object an item is: passive evidence, a preference, a policy, a reusable skill, or an executable commitment. It matters because the consequences of misusing an item scale with its type. The procedural-memory lineage, in which an agent accumulates reusable skills it can later invoke, makes this concrete: Voyager builds an open-ended library of executable skills the agent calls as actions \citep{wang2023voyageropenendedembodiedagent}, and a poisoned skill is far more dangerous than a poisoned fact because invoking it has tool consequences in the world. A reflective lesson degrades gracefully if it is wrong; an executable commitment does not. Actionability therefore sets the bar for the others: the higher an item's actionability, the stronger the authority, provenance, and recoverability guarantees it should require before it is allowed to act. The gap is that current systems rarely type their state by actionability at all, so an advisory note and a callable side-effecting procedure are stored, retrieved, and trusted by the same undifferentiated machinery.

These six axes are the analytic spine of the survey. Every mechanism in later sections can be read as making some axes explicit and leaving others implicit, and every failure mode as the violation of an invariant attached to one or more axes. In our coded corpus, coverage is skewed: mutability and provenance appear more often than authority and recoverability, so the literature attends most to axes whose failures are internal and revisable and less to those whose failures license an action that cannot be taken back.

\subsection{The cognitive-architecture lineage}\label{subsec:lineage}

The axes inherit their structure from the cognitive-architecture tradition, which spent decades formalizing many of the distinctions, memory types, retrieval dynamics, and consolidation processes that LLM agents are now rediscovering empirically. Naming this lineage is not antiquarianism: it explains why the axes are the right ones and why the failures the survey documents were predictable.

The foundational distinction is Tulving's separation of episodic memory, memory for events bound to a time and place, from semantic memory, the store of decontextualized facts \citep{tulving1972episodic}. This underwrites both the provenance axis, since the move from episodic to semantic is precisely the loss of source provenance must guard against, and the state-type taxonomy the survey uses throughout. The complementary-learning-systems theory of McClelland and colleagues then explained why brains keep two stores changing at different rates, a fast hippocampal store that captures specifics and a slow neocortical store that consolidates regularities, arguing formally that interleaving through replay prevents catastrophic interference \citep{mcclelland1995complementary}, an argument later updated to incorporate generative rehearsal and schema-driven fast consolidation \citep{kumaran2016complementary} and grounded empirically in offline hippocampal replay during sleep \citep{wilson1994reactivation}. This is the direct ancestor of the mutability axis and the plasticity-stability tension that organizes the survey's treatment of continual adaptation. Baddeley's multi-component working-memory model, with its episodic buffer binding integrated representations for the current task, supplies the volatile-versus-durable contrast the scope axis formalizes \citep{baddeley2000episodic}, and the quantitative study of forgetting, the Ebbinghaus curve and its modern fitted decay functions, gives the mutability and recoverability axes their measurable shape \citep{murre2015replication}.

These theories were assembled into running architectures. Soar organizes cognition around a decision cycle over explicitly separated procedural, episodic, and semantic memories, treating memory-type isolation as an architectural property enforced by design \citep{laird2012soar}. ACT-R makes the retrieval and action dynamics quantitative, with declarative chunks whose activation decays over time and whose retrieval is probabilistic, and procedural production rules selected by a utility calculus \citep{anderson2004actr}; this is the source of both the activation-decay intuition behind mutability and the possession-versus-licensed-use intuition behind authority. A bridge to LLM agents is CoALA, which recasts the cognitive-architecture frame for language agents, organizing them by memory modules (working, episodic, semantic, procedural) and a structured action space spanning internal and external actions \citep{sumers2024cognitivearchitectureslanguageagents}. CoALA gives a principled vocabulary for much of the modern memory literature without itself enforcing the governance boundaries this survey studies.

Yet the comparison with the classical architectures is also where the lineage exposes its gap. Soar and ACT-R assumed memory-type boundaries were enforced by the architecture, that a procedural rule could not silently corrupt an episodic trace, because the architecture was a closed, hand-built system with hard type boundaries. LLM agents have no such runtime enforcement: their memory types are textual entries in a store, and a procedural skill, a semantic fact, and an episodic trace are often the same kind of object distinguished only by a tag, so the isolation Soar guaranteed by construction must be re-established by governance. CoALA inherits this gap, organizing memory modules as a clean taxonomy but treating them as functionally isolated black boxes, without authority hierarchies that say which module may license an action, without scope gates that prevent one module from corrupting another, and without recoverability handles that allow a corrupted module's effects to be undone. The classical architectures bequeathed the axes but not their enforcement, because enforcement was free in a closed symbolic system and is expensive in an open generative one. The cognitive lineage thus supplies the survey both its vocabulary and its open problem: much of the literature we survey reconstructs the memory \emph{types} of cognitive architecture without reconstructing the governance that made those types safe to use.

\subsection{Inclusions, exclusions, and why personalization and trigger state are state types}\label{subsec:include-exclude}

We close by stating the definition's edges, since a definition that includes everything organizes nothing. It \emph{includes} any system that conditions current behavior on durable cross-session state it accumulates, can act on that state, and can in principle be held to temporal accountability for it. This includes systems the existing literature scatters across separate communities: conversational-memory agents, OS-style memory managers, skill-learning and reflection systems, multi-agent shared-memory platforms, personalized assistants, and proactive or trigger-driven agents. They are unified by a structural commitment to authoritative retained state, rather than by a shared mechanism or application. The definition \emph{excludes} pure architecture or pretraining work that adds no persistent-state operation, episodic agents that reset between tasks, long-context models that consume but do not curate, and read-only RAG over a fixed corpus (Section~\ref{subsec:boundaries}): each lacks the governed accumulation the definition turns on.

The most important boundary call concerns personalization and trigger behavior, because both are sometimes proposed as the defining feature of an always-on agent, and treating them so is a category error this survey avoids. Personalization is a use of persistent state: a personalized agent is one whose retained state includes a user model, and personalization quality is a function of how well that state is written, scoped, and kept current. The personalization benchmarks reviewed above, preference-following \citep{zhao2025prefeval}, persona tracking \citep{jiang2025personamem}, profile-conditioned generation \citep{salemi-etal-2024-lamp}, and the surveys organizing the area \citep{xu2026personalizedllmpoweredagentsfoundations}, all measure capabilities expressed over user state, and measure them well, but they rarely test the governance properties that make the state safe. A useful stress case is authority lapse: a user states a preference, later retracts it, and a personalized agent that still scores well on preference-following may act on the retracted preference, because following a stated preference and respecting a revoked one are different obligations, and only the first is scored. The same holds for trigger and proactivity state, when to interrupt, surface help, or act unprompted \citep{lu2024proactiveagent,he2024madialbench}: a trigger condition is a stored item with authority, scope, and provenance like any other, and proactivity that fires on a stale or out-of-scope trigger is an authority or scope failure, not a personalization success.

We therefore classify personalization, user models, and trigger conditions as persistent-state \emph{types}, taking their place alongside episodic traces, semantic facts, preferences, procedural skills, workflow ledgers, policy and permission state, provenance records, shared and social memory, task ledgers, and tool and credential state in the state-type taxonomy the survey develops next. This keeps the definition faithful to the thesis. If we defined always-on agents by personalization, we would have written a personalization survey with a memory chapter; if by memory, the relabeled memory survey the field already has several of. Defining them as persistent-state systems, and demoting personalization and triggering to state types governed by the same six axes and lifecycle invariants as every other type, keeps the unit of analysis on the object that actually distinguishes always-on agents from their neighbors: not what they remember, and not whom they serve, but the state they make authoritative and whether they can govern it. Later parts answer that governance question at the level of substrates, lifecycle, adaptation, and evaluation.
